% SCP10, source lines 1818--1915. Untwisted tori with coarse periods at least three.
% This is a scope-restricted physical fixed-point theorem, not the unrestricted source label.

\begin{theorem}[Isometric equivalence with the canonical physical support]%
    \label{thm:peps_g_isometric_canonical_support}
    \lean{TNLean.PEPS.IsGIsometric.exists_canonicalSupportEquiv}
    \leanok
    \uses{thm:peps_g_isometric_support_coordinates}
    Let $T$ be $G$-isometric for a unitary representation $\rho$ of a finite
    group, and put $P_\rho=|G|^{-1}\sum_g\rho(g)$. For some $c>0$, the map
    $J=c^{-1/2}T^*$ is an isometric equivalence
    $\operatorname{ran}T\cong\operatorname{ran}P_\rho$, with
    $JT=\sqrt c\,P_\rho$. No surjectivity onto either ambient coordinate
    space is required. This is the support identification used in
    \cite[Observations~6.4 and~6.6]{Schuch2010PEPS}.
\end{theorem}

\begin{proof}\leanok
    The normalized adjoint is already isometric on $\operatorname{ran}T$.
    Its image is contained in $\operatorname{ran}P_\rho$ because
    $JTu=\sqrt c\,P_\rho u$. Conversely, $P_\rho u=JT(c^{-1/2}u)$,
    so it maps onto the canonical support.
\end{proof}

\begin{definition}[Paired boundary labels as regular bundles]%
    \label{def:peps_two_by_two_bundle_coordinates}
    \lean{TNLean.PEPS.pairUnitBundleEquiv,
        TNLean.PEPS.twoByTwoBundledGraphSite,
        TNLean.PEPS.twoByTwoIncidentBoundaryEquiv}
    \leanok
    \uses{def:peps_two_by_two_graph_site}
    Identify $G\times G$ with $G\times G^{\{*\}}$ by
    $(a,b)\mapsto(a,(*\mapsto b))$. Express the actual paired coarse tensor
    in these incident bundle coordinates, without changing its coefficients
    or physical alphabet $P^4$. Reading the four incident edges and the two
    labels on each edge identifies this boundary with all eight separate
    fine boundary legs.
\end{definition}

\begin{theorem}[Exact state and symmetry in bundle coordinates]%
    \label{thm:peps_two_by_two_bundle_coordinates}
    \lean{TNLean.PEPS.graphBondNetwork_twoByTwoBundledGraphSite,
        TNLean.PEPS.torusBondNetwork_eq_twoByTwoBundledGraphSite,
        TNLean.PEPS.regularSiteMap_twoByTwoBundledGraphSite,
        TNLean.PEPS.twoByTwoIncidentBoundaryEquiv_smul,
        TNLean.PEPS.twoByTwoIncidentBoundaryEquiv_rep}
    \leanok
    \uses{def:peps_two_by_two_bundle_coordinates,
        thm:peps_two_by_two_graph_contraction}
    Suppose both coarse periods are at least three. The fine state after
    physical grouping is exactly the actual bundled graph state.
    The eight-label boundary bijection intertwines simultaneous left
    translation and leaves the physical block map unchanged. Thus the
    bundled incident representation is precisely the eight-leg regular
    representation in different coordinates.
\end{theorem}

\begin{proof}\leanok
    Reindex every edge label by the displayed two-label bijection.
    Each boundary coordinate still has the same group value, so translation
    commutes with the bijection. Reindexing the local virtual sum gives the
    equality of physical maps.
\end{proof}

\begin{theorem}[Identification of the original coarse tensor]%
    \label{thm:peps_two_by_two_original_canonical}
    \lean{TNLean.PEPS.graphIncidentRepresentation_leftRegularMatrix,
        TNLean.PEPS.IsGIsometric.isGIsometric_torusIncidentFamily,
        TNLean.PEPS.IsGIsometric.exists_torusCanonicalSupportEquiv}
    \leanok
    \uses{thm:peps_g_isometry_coordinate_equiv,
        thm:peps_g_isometric_canonical_support}
    On any finite simple graph, the oriented incident action of the
    left-regular matrices equals simultaneous left translation on incident
    labels. On a square torus of periods at least three, let $A$ be any
    four-leg regular $G$-isometric tensor. Its incident-coordinate physical
    map $T_A$ is $G$-isometric. At each coarse vertex there are $c>0$ and an
    isometric equivalence between its physical support and the support of
    the canonical regular averaging tensor $T_0$, acting by
    $J=c^{-1/2}T_A^*$ and satisfying $JT_A=\sqrt c\,T_0$.
    This identifies the coarse tensor with the original tensor on physical
    support, with its scalar normalization retained.
\end{theorem}

\begin{proof}\leanok
    The transposed inverse regular matrix on an outgoing leg has the same
    entries $\delta_{a,gb}$ as the incoming matrix. The four-leg incident
    bijection therefore intertwines both representations and preserves
    their coordinate inner products. Apply the canonical support
    equivalence to the transported original tensor.
\end{proof}

\begin{theorem}[G-isometry of actual bundled coarse sites]%
    \label{thm:peps_two_by_two_bundled_isometry}
    \lean{TNLean.PEPS.isGIsometric_twoByTwoBundledGraphSite}
    \leanok
    \uses{thm:peps_two_by_two_isometry,
        thm:peps_two_by_two_bundle_coordinates,
        thm:peps_g_isometry_coordinate_equiv}
    If every fine site is regular $G$-isometric, then every actual bundled
    coarse site is $G$-isometric for its bundled incident representation.
    There is no independent hypothesis about the blocked Gram matrix.
\end{theorem}

\begin{proof}\leanok
    The explicit four-site Gram calculation gives G-isometry on all eight
    separate boundary labels. Transport it through the actual boundary
    coordinate equivalence.
\end{proof}

\begin{definition}[The physical support of the fine state]%
    \label{def:peps_two_by_two_fine_support}
    \lean{TNLean.PEPS.twoByTwoPhysicalRegroupingIsometry,
        TNLean.PEPS.twoByTwoPhysicalRegroupingIsometry_apply,
        TNLean.PEPS.twoByTwoFineState,
        TNLean.PEPS.twoByTwoFinePhysicalSupport,
        TNLean.PEPS.twoByTwoPhysicalSupportRegrouping}
    \leanok
    \uses{def:peps_two_by_two_physical_grouping,
        def:peps_two_by_two_bundle_coordinates}
    Let $\Psi_a$ be the original fine bond-indexed torus state. Let $W$
    be its full physical regrouping unitary, and let $T_B=\bigotimes_vB_v$
    be the product map of the actual bundled coarse tensors. Define
    $S_a=W^{-1}(\operatorname{ran}T_B)$. Restricting $W$ gives an isometric
    equivalence $S_a\cong\operatorname{ran}T_B$. All physical coordinates
    remain in their original ambient spaces.
\end{definition}

\begin{theorem}[Derived support membership of the actual fine state]%
    \label{thm:peps_two_by_two_fine_support}
    \lean{TNLean.PEPS.twoByTwoPhysicalRegroupingIsometry_fineState,
        TNLean.PEPS.twoByTwoFineState_mem_physicalSupport}
    \leanok
    \uses{def:peps_two_by_two_fine_support,
        thm:peps_two_by_two_bundle_coordinates,
        thm:peps_graph_physical_support_transport}
    The physical regrouping satisfies $W\Psi_a=\Psi_B$, and the actual
    fine state belongs to $S_a$. These statements require no symmetry or
    supplied support hypothesis.
\end{theorem}

\begin{proof}\leanok
    The geometric bond bijection gives $W\Psi_a=\Psi_B$.
    The graph contraction is the product of the local block maps applied
    to its virtual bond vector, hence $\Psi_B\in\operatorname{ran}T_B$.
\end{proof}

\begin{theorem}[Physical four-site renormalization on a regular torus]%
    \label{thm:peps_two_by_two_physical_renormalization}
    \lean{TNLean.PEPS.exists_regularTwoByTwoPhysicalSupportIsometry}
    \leanok
    \uses{thm:peps_two_by_two_bundled_isometry,
        thm:peps_two_by_two_fine_support,
        thm:peps_regular_graph_physical_blocking}
    Let $G$ be any finite group, and let the two coarse periods be at least
    three. Place arbitrary regular $G$-isometric fine tensors $a_z$ on the
    doubled torus, with any common finite physical alphabet. There are
    positive block factors $c_v$ and a Hilbert-space isometry $I$ on the
    derived support $S_a$ such that
    \begin{align}
        I\Psi_a=
        \left(\prod_v\sqrt{c_v}\right)
        (\sqrt{|G|})^{|E|}\,\Psi_0\otimes\Omega,
        \qquad \|\Omega\|^2=1.
    \end{align}
    Here $\Psi_0$ is the actual canonical regular PEPS on the coarse
    graph, and $\Omega$ is one normalized maximally entangled group-label
    pair on each coarse edge. Explicitly,
    $I=U(\bigotimes_vc_v^{-1/2}B_v^*)W$ on $S_a$, where $U$ takes the
    relative group coordinate on every paired physical half-edge.
    The four-site Gram calculation derives block factors
    $|G|\prod_{i=0}^3c_{v,i}$ from positive fine-site factors.
    For a common original tensor $A$, the coarse canonical tensor is
    identified with $A$ by the preceding physical-support equivalence.
    This is the untwisted periodic specialization of
    \cite[Observations~6.5--6.6]{Schuch2010PEPS}.
\end{theorem}

\begin{proof}\leanok
    Apply physical support accessibility and regular bundle separation to
    the actual coarse graph tensors, whose G-isometry has been derived
    from the fine tensors. Compose its support isometry with the restricted
    physical regrouping $W$. The exact state equality supplies the required
    state argument, and its support membership was derived separately.
    Each unnormalized surplus equality vector is $\sqrt{|G|}$ times its
    normalized Bell vector, producing the displayed factor on all edges.
\end{proof}
